%% ============================================================================
\section{Relation Mining and Gold Weights}
\label{app:estimation}
%% ============================================================================

\paragraph{Pairwise miner.} Prompt~A receives the full response and one claim pair and returns a
discourse sense and coupling, with reasoning \emph{before} the label so that the label's token
log-probabilities measure a decision already made; the type confidence is the per-token geometric
mean probability of the label span, or a self-consistency estimate
\citep{simbauq} on backends without log-probabilities. Prompt~B asks whether, given the coupling,
the source makes the target at least plausible; the answer must begin with \texttt{Yes} or
\texttt{No}, and the strength is the renormalized $P(\texttt{Yes})$ \citep{kadavath}, or the
affirmative fraction over samples. A verbalized \texttt{[p=0.NN]} variant is kept only as a
comparison, since verbalized confidence is weakly calibrated \citep{xiong}. Two precision
mechanisms are added to grounding. An \emph{admissibility filter} restricts answers to the
(sense, coupling) combinations of Table~\ref{tab:couplings} that the corpus uses: the ten senses
of App.~\ref{sec-full-locobench}, excluding \sense{Effect-Cause}, \sense{Condition} and
\sense{Instantiation}, whose direction readers dispute or which fit almost any topical pair. An
\emph{evidence-span requirement} obliges the model to quote, verbatim from the response, the words
that draw the link, and the parser drops the relation if the quotation is not a substring of the
response.

\paragraph{Candidate pairs.} Querying all $n(n-1)$ ordered pairs is quadratic in calls and, without
grounding, yields over-connected graphs. Four policies are available: \textsc{all\_pairs};
\textsc{windowed}, forward pairs with $0<j-i\le w$ ($w=4$); \textsc{gated}, the window plus
long-range forward pairs above an embedding or entity similarity threshold; and
\textsc{bidirectional}, both directions within $|j-i|\le w$. The last exists because a forward-only
policy cannot emit a relation that is both directed and backward in claim order. Every policy is
refined against the response: a pair is \emph{discourse-adjacent} if the two claims share content
above a Jaccard threshold, if their sentences are close, or if the target's sentence opens with a
connective. Out-of-window adjacent pairs are promoted into the candidate set, and in-window
non-adjacent pairs are demoted out of it.

\paragraph{Prompt revisions and failure accounting.} A first-generation Prompt~A over-asserted,
answering ``related'' on most candidate pairs. The revision rebalances the few-shot exemplars
towards unrelated pairs, uses silence rather than explicit disavowal as the negative signal, asks
the discriminating question (does the text draw the link, or are the claims merely near each
other?) first, and names transitivity, closing a chain the text laid out pairwise, as a
non-relation. Mining issues one Prompt~A call per candidate pair and one Prompt~B call per
resulting edge. A failed call returns no parse, and a parser that read ``no parse'' as ``no
relation'' would turn an infrastructure failure into a sparser, and therefore more coherent-looking,
graph; call errors are therefore counted separately from genuine \emph{none} answers.

Both prompts teach the compile map of Table~\ref{tab:couplings}: \sense{Refutation} is the only
sense that compiles to \con{}, and \sense{Contrast} and \sense{Concession} compile to \emph{none}.
In the experiments the pairwise miner uses the response-grounded Prompt~A with the evidence-span
requirement and the admissibility filter, type confidences from the label tokens'
log-probabilities and strengths from the \texttt{Yes}/\texttt{No} surrogate, over the window
policy with $w=4$.

\paragraph{Graph-level miner.} The prompt hands the model the response and the claims numbered from
$1$, and asks for a JSON list of every pair the response relates, each with a sense from the ten of
the corpus, its coupling and a strength band (strong, moderate, weak). Its rules: emit a relation
only if the response itself draws it; orient every entailment from the supporting claim to the
supported one; separate \sense{Refutation} from \sense{Contrast} by whether the response asserts
incompatibility; and relate a pair only once. Each response is sampled five times (gpt-oss-120b at
temperature $0.7$, GPT-6 at the provider default), and relations are aggregated by strict-majority
vote as in \S\ref{sec:mining}. The committee that admitted the corpus never saw this prompt: it
confirmed relations pair by pair with a classification prompt (App.~\ref{sec-full-locobench}). The
prompt follows verbatim.
\promptfile{prompts/PROMPT_GRAPH_V2.txt}

\paragraph{Gold arm.} Each gold relation's weight is the midpoint of its labeled strength band
(strong $[0.85,1]$, moderate $[0.60,0.84]$, weak $[0.35,0.59]$). The corpus's factuality labels give
the committed world ($\xs_i=1$ for claims labeled factual, $0$ otherwise), and each response's
stance on each claim gives the world of assertion polarity. The score's MRF uses flat priors. No miner is a member of the committee that admitted
the corpus's items.

\paragraph{Pairwise miner.} The pairwise miner asks one question per candidate pair within a window
of four claims, and reads $c_r$ from the label-token probabilities and $g_r$ from a
\texttt{Yes}/\texttt{No} judgment (\S\ref{sec:mining}). Its prompt follows verbatim.
\promptfile{prompts/PROMPT_PAIR_SENSE_V2.txt}

\paragraph{Judges.} Both judges see the response only, not its claims or verdicts. G-Eval uses the
SummEval coherence rubric as published; the direct judge asks for a 1--5 rating of logical
coherence. Each response is rated five times, with the gateway's response cache bypassed. The
GPT-6 judges sample an integer rating; the gpt-oss-120b judges report the rating's expectation under
the rating token's probabilities, which is why they almost never tie. The prompts follow verbatim.
\promptfile{prompts/JUDGE_GEVAL.txt}
\promptfile{prompts/JUDGE_DIRECT.txt}

\paragraph{Fairer judges.} Three further GPT-6 judges answer review concerns about the comparison:
an absolute 1--5 rating given the claims marked TRUE or FALSE and our definition of logical
coherence (three runs averaged); and a pairwise judge, asked which of two rungs is more coherent,
in both orders, with and without the claims, verdicts and definition. A pair counts as a rise when
the net preference over the two orders favors the declared rung, and as a tie when it is zero. The
gateway cache was bypassed. The prompts follow verbatim; the essay grader of App.~\ref{app:asap}
uses the last one.
\promptfile{prompts/JUDGE_VERDICT_ABS.txt}
\promptfile{prompts/JUDGE_PAIRWISE_VERDICTS.txt}
\promptfile{prompts/JUDGE_PAIRWISE.txt}
\promptfile{prompts/GRADER_ASAP_RUBRIC.txt}
The 0--100 essay grader of App.~\ref{app:asap} uses the same rubric with this final instruction:
\promptfile{prompts/GRADER_ASAP_RUBRIC100.txt}

\paragraph{Other baselines.} DiscoScore uses \texttt{bert-base-uncased}. NLI contradiction and
ROSCOE self-consistency use gpt-oss-120b as the NLI model and read only the claims, not the
response; the claims are identical across a ladder's rungs, so both tie on every pair by
construction and are omitted from Figure~\ref{fig:ladders}. The signed weight sum and the signed
count (each relation $+p_r$ or $+1$ if it is a support, $-p_r$ or $-1$ if it is a conflict) read
the same graphs and committed worlds as $\LCS$.

%% ============================================================================
\section{Discourse Senses, by Example}
\label{app:senses}
%% ============================================================================

Table~\ref{tab:senseex} gives one response excerpt per sense of Table~\ref{tab:couplings}, grouped
by the coupling each compiles to. Every excerpt contains both claims (\uline{$A$ underlined},
\textbf{$B$ bold}); the \emph{italic} words draw the link, and they are what an annotator, human or
model, must recognize. The \sense{Refutation}, \sense{Alternative} and \sense{Disjunction} rows
describe one two-part weld failure three ways, which differ only in exhaustiveness, that is, in
which same-value worlds the response leaves available; the \sense{Contrast} row shows the same
material without an asserted incompatibility. Table~\ref{tab:sensemap} gives the closest PDTB~3.0
and RST label for each sense. The corpus of \S\ref{sec:experiments} uses every sense except
\sense{Effect-Cause}, \sense{Condition} and \sense{Instantiation}: the first and last run in a
direction that readers dispute, and \sense{Condition} fits almost any pair of claims on one
topic.

\begin{table*}[t]
\centering\footnotesize
\setlength{\tabcolsep}{4pt}
\renewcommand{\arraystretch}{1.12}
\begin{tabularx}{\textwidth}{@{}llX@{}}
\toprule
Sense & Coupling & Response excerpt, and why this sense rather than a neighbor \\
\midrule
\sense{Cause-Effect} & \ent
  & ``\uline{The company shipped a batch it knew to be out of tolerance}, and \emph{as a direct
    result} \textbf{its share price fell $15\%$ that quarter}.'' The response presents $A$ as
    bringing $B$ about. \\
\sense{Effect-Cause} & \ent
  & ``\uline{The entire fleet was grounded on 14 March}. The order came \emph{because}
    \textbf{inspectors had found a cracked wing spar on two airframes}.'' The direction of
    explanation is reversed; the direction of inferential support is not. \\
\sense{Evidence} & \ent
  & ``\uline{The maintenance log records no inspection between January and June}, \emph{which
    shows that} \textbf{the mandatory 500-hour check was skipped}.'' $A$ supports $B$ without
    causing it: a record is evidence, not a mechanism. \\
\sense{Condition} & \ent
  & ``\textbf{A part is cleared for airframe use} \emph{only if} \uline{its alloy meets the
    reference specification}.'' $B$ is asserted contingent on $A$; nearly any dependence can be
    reread as a condition. \\
\sense{Instantiation} & \ent
  & ``\uline{Several suppliers failed the 2023 audit} --- \emph{for instance}, \textbf{the Tucson
    plant failed on two counts}.'' $B$ is one case of $A$'s generalization. \\
\midrule
\sense{Restatement} & \eqv
  & ``\uline{The shipment was certified to specification} --- \emph{in other words},
    \textbf{every requirement in the contract had been met}.'' A marked paraphrase asserts
    identity of content. \\
\midrule
\sense{Refutation} & \con
  & ``\uline{The weld passed the tensile test}, \emph{which rules out} \textbf{the same joint
    being porous along its full length}.'' Asserted incompatibility, \emph{not} exhaustive: if the
    wrong joint was tested, both are false, so the both-false world stays available. \\
\midrule
\sense{Alternative} & \exc
  & ``The two accounts cannot be reconciled: \emph{either} \uline{no one aboard was harmed},
    \emph{or} \textbf{three passengers died}, and \emph{exactly one of them is true}.''
    Exhaustive: the pair can neither both hold nor both fail. \\
\midrule
\sense{Disjunction} & \cnc
  & ``The fracture began at the joint, so \emph{at least one of} \uline{the weld} \emph{and}
    \textbf{the fastener} \emph{must have failed} --- possibly both.'' The dual of
    \sense{Alternative}: only the both-false world is ruled out. \\
\midrule
\sense{Contrast} & none
  & ``\uline{The weld passed the tensile test}, \emph{but} \textbf{the adjacent rivets showed
    corrosion}.'' Both can hold; the \emph{but} marks a contrast of expectations, not an
    incompatibility, so the sense induces no factor. \\
\sense{Concession} & none
  & ``\emph{Although} \uline{the airline argued throughout that pilot error was to blame},
    \textbf{the final report identified a metallurgical defect}.'' Both claims hold (the airline
    did argue it); the tension is between the airline's position and the finding, so no factor. \\
\sense{Precedence} & none
  & ``\uline{The replacement units were installed in March}, \emph{some three months before}
    \textbf{the first in-flight failures were reported in June}.'' The text relates the claims
    and says nothing about either one's truth, so the sense induces no factor. \\
\sense{Succession} & none
  & ``\textbf{The first in-flight failures were reported in June}, \emph{fully three months
    after} \uline{the replacement units had been installed}.'' The mirror of
    \sense{Precedence}. \\
\bottomrule
\end{tabularx}
\caption{One response excerpt per discourse sense of Table~\ref{tab:couplings}.}
\label{tab:senseex}
\end{table*}

\begin{table}[t]
\centering\footnotesize
\setlength{\tabcolsep}{3pt}
\begin{tabular}{@{}lll@{}}
\toprule
Sense & PDTB~3.0 & RST \\
\midrule
\sense{Cause-Effect} & Contingency.Cause.Result & Result \\
\sense{Effect-Cause} & Contingency.Cause.Reason & Cause \\
\sense{Evidence} & Contingency.Cause+Belief & Evidence \\
\sense{Condition} & Contingency.Condition & Condition \\
\sense{Instantiation} & Expansion.Instantiation & Elaboration \\
\sense{Restatement} & Expansion.Equivalence & Restatement \\
\sense{Disjunction} & Expansion.Disjunction & Disjunction \\
\sense{Alternative} & Expansion.Disjunction & Disjunction \\
\sense{Refutation} & --- & Antithesis \\
\sense{Contrast} & Comparison.Contrast & Contrast \\
\sense{Concession} & Comparison.Concession & Concession \\
\sense{Precedence} & Temporal.Asynch.Precedence & Sequence \\
\sense{Succession} & Temporal.Asynch.Succession & Sequence \\
\bottomrule
\end{tabular}
\caption{Closest PDTB~3.0 and RST labels for each sense; \sense{Alternative} is the exclusive
reading of Expansion.Disjunction. PDTB~3.0 has no sense for an asserted incompatibility; \sense{Refutation} is closest to RST's \emph{Antithesis}. The mapping is
many-to-one in places, so senses are annotated directly rather than converted.}
\label{tab:sensemap}
\end{table}

